\section{Results}
\label{sec:results}

\subsection{Reproducing MultiResNetSCA-1}
Table~\ref{tab:heads} compares our reproduction with Table~2 of~\cite{MS23}. The four head
accuracies agree to within two points ($\alpha$ 99.94 vs.\ 99.2, $\beta$ 19.42 vs.\ 21.1, $C$ 1.51
vs.\ 1.6, $\pi$ 88.77 vs.\ 88.9\,\%), so the residual-network rows below are a faithful same-pipeline
reference, not a weakened baseline.

\begin{table}[t]
\centering
\caption{Test accuracy of each head on 5,000 attack traces ($C$ and $\pi$: mean over the 16 heads;
$C$ gain: mean log-likelihood of the true label minus that of a uniform guess).}
\label{tab:heads}
\resizebox{\linewidth}{!}{% generated by paper/tches/scripts/make_tables.py -- do not edit by hand
\begin{tabular}{lrrrrr}
\toprule
Model & $\alpha$ (\%) & $\beta$ (\%) & $C$ (\%) & $C$ gain (nats) & $\pi$ (\%) \\
\midrule
MS23 MultiResNetSCA-1 (Table 2 of~\cite{MS23}) & 99.2 & 21.1 & 1.6 & -- & 88.9 \\
MultiResNet (reprod.), best & 99.94 & 19.42 & 1.51 & +0.344 & 88.77 \\
MultiResNet (reprod.) s2, best & 100.00 & 20.40 & 1.43 & +0.355 & 88.26 \\
MultiResNet + $w_C$, $\mathcal{L}_{\mathrm{dec}}$, best & 100.00 & 20.06 & 1.40 & +0.324 & 88.49 \\
Compact (plain), best & 99.78 & 22.18 & 1.50 & +0.398 & 90.98 \\
Compact s1, best & 99.94 & 23.28 & 1.53 & +0.390 & 91.08 \\
Compact s2, best & 99.80 & 23.44 & 1.47 & +0.383 & 90.64 \\
GPAM, paper LR & 51.98 & 2.76 & 0.38 & -0.002 & 7.79 \\
GPAM, LR $10^{-3}$ & 99.88 & 5.72 & 0.93 & +0.139 & 55.05 \\
GPAM, LR $10^{-3}$ + sub\_avg & 99.68 & 8.76 & 1.00 & +0.243 & 54.46 \\
\bottomrule
\end{tabular}
}
\end{table}

\subsection{Full-key recovery}
Table~\ref{tab:main} reports every model on the same 5,000 traces with the decoder of
Eq.~\eqref{eq:decode}. Under the metric of~\cite{MS23} (GE$<1$ on all bytes, 500 attacks),
CompactMultiResNet needs 60 traces with both seeds -- the figure reported for MultiResNetSCA-1
in~\cite{MS23} -- and our 137.4M reproduction needs 75, also with both seeds. We therefore claim parity with the published
attack at 8.8$\times$ fewer parameters and a 1.25$\times$ improvement over a same-pipeline
reproduction, not a new trace record.

\begin{table}[t]
\centering
\caption{Full-key recovery on 5,000 ASCADv2 attack traces, 400k profiling traces, independent
permutation heads. Metrics as defined in Section~\ref{sec:metrics}; GE columns over 500 attacks.}
\label{tab:main}
\resizebox{\linewidth}{!}{% generated by paper/tches/scripts/make_tables.py -- do not edit by hand
\begin{tabular}{llrrrrrr}
\toprule
Model & Params (M) & 16/16 & median $z$ & GE$<1$ & GE$<0.5$ & GE$=0$ & MTD \\
\midrule
MS23 MultiResNetSCA-1~\cite{MS23} & 137.4 & yes & -- & $\approx$60 & -- & -- & -- \\
MultiResNet (reprod.), best & 137.4 & yes & +14.21 & 75 & 90 & 235 & 47.4 \\
MultiResNet (reprod.), last & 137.4 & yes & +13.25 & 160 & 190 & 450 & 84.0 \\
MultiResNet (reprod.) s2, best & 137.4 & yes & +14.28 & 75 & 85 & 185 & 43.0 \\
MultiResNet (reprod.) s2, last & 137.4 & yes & +13.43 & 135 & 150 & 330 & 76.2 \\
MultiResNet + $w_C$, $\mathcal{L}_{\mathrm{dec}}$, best & 137.4 & yes & +14.23 & 80 & 90 & 235 & 47.0 \\
MultiResNet + $w_C$, $\mathcal{L}_{\mathrm{dec}}$, last & 137.4 & yes & +13.49 & 150 & 175 & 405 & 86.1 \\
Compact (plain), best & 15.65 & yes & +14.40 & 60 & 70 & 150 & 35.9 \\
Compact (plain), last & 15.65 & yes & +14.22 & 70 & 80 & 200 & 41.6 \\
Compact s1, best & 15.65 & yes & +14.44 & 60 & 70 & 165 & 35.8 \\
Compact s1, last & 15.65 & yes & +14.37 & 70 & 80 & 190 & 40.9 \\
Compact s2, best & 15.65 & yes & +14.41 & 60 & 70 & 170 & 37.0 \\
Compact s2, last & 15.65 & yes & +14.28 & 75 & 80 & 175 & 41.7 \\
GPAM, paper LR & 6.11 & 0/16 & -0.24 & -- & -- & -- & -- \\
GPAM, LR $10^{-3}$ & 6.11 & yes & +8.65 & 1145 & 1295 & 2475 & 599.2 \\
GPAM, LR $10^{-3}$ + sub\_avg & 6.11 & yes & +10.28 & 550 & 625 & 1270 & 228.8 \\
\bottomrule
\end{tabular}
}
\end{table}

\paragraph{Architecture or losses?} CompactMultiResNet is trained with $w_C=0.5$ and $\Ldec$. Giving
the 137.4M network the same two losses does not help it (GE$<1$ 80 vs.\ 75 traces, median $z$
+14.23 vs.\ +14.21); conversely, CompactMultiResNet trained \emph{without} them needs the same 60
traces (GE$=0$ 150). At 400k profiling traces the gap to the 137.4M network is therefore due to the
architecture, not to the two losses (one seed per ablation row).

\subsection{Robustness to checkpoint choice}
With both seeds the 137.4M network reaches its lowest validation loss at epoch 4 and then overfits
(seed 1: summed validation cross-entropy of the 34 heads 89.7 $\to$ 104.0 nats at epoch 15). Its last checkpoints need 160 and 135 traces (GE$<1$), and
their $C$ heads score \emph{below} a uniform guess on their own labels ($-0.48$ and $-0.40$ nats). With the same losses as
CompactMultiResNet it degrades in the same way. CompactMultiResNet's last checkpoints need 70 and 75
traces (plain: 70), and their $C$ gains stay positive (+0.31 for seed 1). This is consistent with the
excess head capacity being spent on memorising the profiling set.


\subsection{A long-range transformer at equal budget}
GPAM was designed for raw 1M-sample traces~\cite{GPAM}; we evaluate it on the 15k extract because
that is the input every other row uses. With the published target learning rate it stays near chance
for 30 epochs and ends with $\alpha$ at 52\% and $C$ at chance: no byte is recovered. With a learning
rate of $10^{-3}$ and CompactMultiResNet's exposure it recovers all 16 bytes but needs 1,145 traces
(GE$<1$), about 19$\times$ more. With the residual networks' mean-trace subtraction, i.e.\ the
same preprocessing as every other row, it needs 550 (9$\times$ more; median $z$ +10.28 vs.\
+8.65). Its validation loss was
still decreasing at the end of training, so these are budget-matched rather than converged numbers.
On raw traces~\cite{GPAM} reports about 80 traces; our reproduction of that setting needed 600
\todo{cite T-019 protocol in appendix}.
